% !TEX root = ../main.tex
\chapter{从模形式到自守表示}\label{ch:07}

\begin{center}
Stephen S. Kudla
\end{center}

本章勾勒一条过渡：从权为 $k$、级为 $N$ 的经典全纯模形式 $f$，走向 $\GL_2(\A)$ 的自守表示 $\pi=\pi(f)$。

在远为简单的 $\GL_1$ 情形，类似的过渡就是 Tate 论文一章所述的对应：模 $N$ 的经典 Dirichlet 特征 $\chi_N$，对应连续特征 $\omega:\GL_1(\mathbb{Q})\backslash\GL_1(\A)\to\mathbb{C}^\times$。$\GL_1$ 情形的两个特点会延续到 $\GL_2$ 情形。第一，特征 $\omega$ 是局部特征之积 $\omega=\otimes_v\omega_v$；例如 Tate 论文中的分析，正是从局部成分 $\omega_v$ 出发建立 $\omega$ 的 $L$ 函数的整体理论。类似地，$\GL_2(\A)$ 的自守表示 $\pi$ 会有分解
\[
\pi=\otimes_v\pi_v,
\]
其中每个 $\pi_v$ 是群 $\GL_2(\mathbb{Q}_v)$ 的表示。第二，同一个 $\omega$ 对应许多经典 Dirichlet 特征，其中恰有一个本原特征。类似地，同一个 $\pi$ 会对应无穷多个 $f$（我们允许级 $N$ 任意变化），但其中恰有一个本原形式，即归一化的新形式。

当然，主要差别在于：当 $G=\GL_2$ 时，表示 $\pi$ 与各局部成分 $\pi_v$ 都是无限维的，因此涉及的表示论要沉重得多，这就是非交换调和分析。此外，$\pi$ 将实现为阿代尔群 $G(\A)$ 通过右平移在 $G(\mathbb{Q})\backslash G(\A)$ 上某个函数空间上的“作用”。稍后我们会简要讨论一般数域的情形。

与 Tate 论文一样，本章的主题在许多地方都有很好、很详细的叙述。本主题的基础性工作是 Jacquet 与 Langlands 的书 [11]，它于 1970 年问世，比 Tate 论文晚二十年。在众多评注与导引中，我们只举 Bump [2]、Gelbart [8]、Godement [10] 的书，以及 Diamond 与 Im 的文章 [7]；其中任何一部都可查到更多细节。若要了解 Langlands 纲领的更广阔图景，读者可参阅 Gelbart [9] 与 Knapp [12] 的导引性文章，以及 Cogdell 的综述 [4, 5, 6]。

我们构造的输入，是一个经典全纯尖点形式 $f\in S_k(\Gamma_0(N),\chi_N)$，其中 $\chi_N$ 是模 $N$ 的经典 Dirichlet 特征。于是 $f$ 是上半平面 $\mathbb{H}=\{\tau=u+iv\}$ 上的全纯函数，并且对所有
\[
\gamma\in\Gamma_0(N)=\left\{\begin{pmatrix}a&b\\c&d\end{pmatrix}\in\mathrm{SL}_2(\mathbb{Z})\ \middle|\ c\equiv 0\bmod(N)\right\}
\]
都满足
\begin{equation*}\tag{1.1}
f\!\left(\frac{a\tau+b}{c\tau+d}\right)=\chi_N(d)\,(c\tau+d)^k\,f(\tau).
\end{equation*}
$f$ 的 $q$ 展开是
\begin{equation*}\tag{1.2}
f(\tau)=\sum_{n=1}^{\infty}a_nq^n,\qquad q=e(\tau).
\end{equation*}
注意：要使非零的 $f$ 存在，必须成立相容性条件 $\chi_N(-1)=(-1)^k$。

\section{与经典尖点形式相联系的 $G(\mathbb{Q})\backslash G(\A)$ 上的函数}

取 $G=\GL_2$，视其为 $\mathbb{Q}$ 上的代数群，并令
\[
G_v=G(\mathbb{Q}_v)=\begin{cases}
\GL_2(\mathbb{Q}_p), & \text{若 }\mathbb{Q}_v=\mathbb{Q}_p,\\
\GL_2(\mathbb{R}), & \text{若 }\mathbb{Q}_v=\mathbb{R}.
\end{cases}
\]
$G$ 的阿代尔点群是
\[
G(\A)=G(\mathbb{R})G(\A_f),\qquad
G(\A_f)={\prod_{p}}'G(\mathbb{Q}_p),
\]
它是关于子群 $K_p=\GL_2(\mathbb{Z}_p)\subset\GL_2(\mathbb{Q}_p)=G_p$ 的限制乘积。于是 $G(\A)$ 与 $G(\A_f)$ 都是局部紧拓扑群，
\[
K=\prod_p K_p\subset G(\A_f)
\]
是紧开子群，而 $G(\mathbb{Q})$ 在对角嵌入
\[
G(\mathbb{Q})\hookrightarrow G(\A),\qquad \gamma\mapsto(\dots,\gamma,\dots)
\]
下的像是 $G(\A)$ 的离散子群。

对群 $\mathrm{SL}_2$ 作类似构造，得到局部紧群 $\mathrm{SL}_2(\A)$、$\mathrm{SL}_2(\A_f)$，以及离散子群 $\mathrm{SL}_2(\mathbb{Q})\hookrightarrow\mathrm{SL}_2(\A)$。强逼近定理\footnote{此时情形的一个初等证明见 [19, p.\ 145]。}断言：子群 $\mathrm{SL}_2(\mathbb{Q})\mathrm{SL}_2(\mathbb{R})$ 在 $\mathrm{SL}_2(\A)$ 中稠密；因此对任意开子群 $U\subset\mathrm{SL}_2(\A_f)$，有
\[
\mathrm{SL}_2(\A)=\mathrm{SL}_2(\mathbb{Q})\,\mathrm{SL}_2(\mathbb{R})\,U,\qquad g=\gamma\,g_\infty\,u.
\]
\begin{annote}
强逼近定理对 $\mathrm{SL}_2$ 成立而对 $\GL_2$ 不成立，这一点是本章若干论证的隐形支点。对 $\GL_2$，$G(\mathbb{Q})G(\mathbb{R})^+K'$ 一般真包含于 $G(\A)$，缺的部分由 $\det$ 控制，大小等于 $\mathbb{Q}$ 的理想类数（对 $\mathbb{Q}$ 平凡，对一般数域则不然）。引理 1.1 之所以仍能断言 $G(\mathbb{Q})\backslash G(\A)/K'\simeq\Gamma'\backslash G(\mathbb{R})^+$，关键在于额外假设了 $\det(K')=\widehat{\mathbb{Z}}^\times$：行列式方向上的逼近问题被 $\mathbb{R}^\times_+$ 与 $\widehat{\mathbb{Z}}^\times$ 的乘积吸收，这里用到了 $\A^\times=\mathbb{Q}^\times\mathbb{R}^\times_+\widehat{\mathbb{Z}}^\times$。若去掉 $\det$ 条件，$\Gamma'$ 会变成若干个 $\Gamma_0(N)$ 的并，Hecke 算子理论须相应修改。

\par\vspace{0.15em}{\raggedleft 【来源：Shimura 上书第 1 章（原书脚注 1 引 [19, p.\ 145]）；Platonov–Rapinchuk《Algebraic Groups and Number Theory》第 7 章（A）】\par}
\end{annote}

利用这一结果以及分解
\[
\A^\times=\mathbb{Q}^\times\mathbb{R}^\times_+\,\widehat{\mathbb{Z}}^\times,
\]
我们得到下述结论。

\medskip\noindent\textbf{引理 1.1.}\quad
设 $K'\subset G(\A_f)$ 是紧开子群，且 $\det(K')=\widehat{\mathbb{Z}}^\times$。则
\[
G(\A)=G(\mathbb{Q})\,G(\mathbb{R})^+\,K'.
\]
特别地，
\[
G(\mathbb{Q})\backslash G(\A)/K'\simeq\Gamma'\backslash G(\mathbb{R})^+,
\]
其中 $G(\mathbb{R})^+=\{g\in G(\mathbb{R})\mid \det(g)>0\}$，而 $\Gamma'=G(\mathbb{Q})\cap G(\mathbb{R})^+K'$，经由到阿基米德成分的投影视为 $G(\mathbb{R})^+$ 的子群。

\medskip
取 $K'=K$ 如上，便得到 $\Gamma'=\mathrm{SL}_2(\mathbb{Z})\subset\GL_2(\mathbb{R})^+$。我们也可以取
\begin{equation*}\tag{1.3}
K'=K_0(N)=\left\{k=\begin{pmatrix}a&b\\c&d\end{pmatrix}\ \middle|\ c\equiv 0\bmod(N)\right\},
\end{equation*}
此时 $\Gamma'=\Gamma_0(N)\subset\GL_2(\mathbb{R})^+$。注意存在满同态\footnote{这其实是指
\[
a=(\dots,a_p,\dots)\mapsto\prod_{p\mid N}\bigl(a_p\bmod(p^{n_p})\bigr)\in\prod_{p\mid N}(\mathbb{Z}/p^{n_p}\mathbb{Z})^\times\simeq(\mathbb{Z}/N\mathbb{Z})^\times,
\]
其中 $n_p=\operatorname{ord}_p(N)$。}
\begin{equation*}\tag{1.4}
K_0(N)\to(\mathbb{Z}/N\mathbb{Z})^\times,\qquad
\begin{pmatrix}a&b\\c&d\end{pmatrix}\mapsto a\bmod(N),
\end{equation*}
于是经典 Dirichlet 特征 $\chi_N$ 定义了 $K_0(N)$ 的一个特征 $\chi$。

群 $G(\mathbb{R})^+$ 通过分式线性变换作用在上半平面上，其中心 $Z_\infty=\{z\cdot\mathbf{1}_2\mid z\in\mathbb{R}^\times\}$ 作用平凡。对 $g\in G(\mathbb{R})^+$ 与 $\tau\in\mathbb{H}$，置
\begin{equation*}\tag{1.5}
j(g,\tau)=\det(g)^{-\frac12}(c\tau+d),\qquad g=\begin{pmatrix}a&b\\c&d\end{pmatrix}.
\end{equation*}
注意当 $z\in Z_\infty$ 时 $j(zg,\tau)=\operatorname{sgn}(z)j(g,\tau)$。又注意若
\begin{equation*}\tag{1.6}
k_\theta=\begin{pmatrix}\cos(\theta)&\sin(\theta)\\-\sin(\theta)&\cos(\theta)\end{pmatrix}\in SO(2)=:K_\infty^+\subset G(\mathbb{R})^+
\end{equation*}
是 $G(\mathbb{R})^+$ 的极大紧子群，则 $Z_\infty K_\infty^+$ 是 $i$ 的稳定化子，并且
\begin{equation*}\tag{1.7}
j(zk_\theta,i)=\operatorname{sgn}(z)\cdot e^{-i\theta}.
\end{equation*}

\begin{annote}
(1.5) 中的 $\det(g)^{-1/2}$ 一项极易被当作装饰而忽略。它以 $\det$ 的平方根归一化自守因子，使链锁性质 $j(g_1g_2,\tau)=j(g_1,g_2\tau)j(g_2,\tau)$ 对 $g_1,g_2\in G(\mathbb{R})^+$ 成立；若只用 $c\tau+d$，当 $\det g\ne1$ 时该性质会多出 $\det(g_1)^{-1/2}$ 的偏差。古典情形 $g=\gamma\in\Gamma_0(N)$ 有 $\det\gamma=1$，故 (1.1) 与 (1.5) 并不冲突；但第 1 节随后要把 $j$ 用到 $\diag(a,a^{-1})$、$k_\theta$ 这类 $\det\ne1$ 的元素上（见 (1.15) 与 (1.7)），此时该因子与 $\operatorname{sgn}(z)$ 一起决定 $\phi$ 在中心方向的变换。另需注意 $\det(g)^{-1/2}$ 只对 $\det g>0$ 取实值，这正是 (1.5) 把定义域限制在 $G(\mathbb{R})^+$ 而非 $G(\mathbb{R})$ 的原因。

\par\vspace{0.15em}{\raggedleft 【来源：Bump《Automorphic Forms and Representations》第 1、2 章；Shimura《Introduction to the Arithmetic Theory of Automorphic Forms》第 1 章（A）】\par}
\end{annote}

有了这些结构背景，我们回到尖点形式 $f\in S_k(\Gamma_0(N),\chi_N)$。对 $g_\infty\in G(\mathbb{R})^+$，令
\begin{equation*}\tag{1.8}
\phi(g_\infty)=\phi_f(g_\infty)=f(g_\infty(i))\,j(g_\infty,i)^{-k}.
\end{equation*}
注意 $\phi$ 是 $G(\mathbb{R})^+$ 上的光滑函数，并且 $f$ 可由它还原：
\begin{equation*}\tag{1.9}
f(\tau)=f(g_\infty(i))=\phi(g_\infty)\,j(g_\infty,i)^{k},
\end{equation*}
其中 $g_\infty$ 是满足 $\tau=g_\infty(i)$ 的任意元素。

\medskip\noindent\textbf{引理 1.2.}\quad
(i) 对 $\gamma\in\Gamma_0(N)$，
\[
\phi(\gamma g_\infty)=\phi(g_\infty)\,\chi_N(d),\qquad \gamma=\begin{pmatrix}a&b\\c&d\end{pmatrix}.
\]
(ii) 对 $zk_\theta\in Z_\infty K_\infty$，
\[
\phi(g_\infty zk_\theta)=\phi(g_\infty)\,\operatorname{sgn}(z)^k\,(e^{i\theta})^k.
\]

\medskip
断言 (i) 由 $f$ 的变换律与自守因子的链锁性质
\[
j(g_1g_2,\tau)=j(g_1,g_2(\tau))\,j(g_2,\tau)
\]
推出。断言 (ii) 由 (1.7) 立得。

最后，利用引理 1.1 的双陪集分解（取 $K'=K_0(N)$），把 $\phi=\phi_f$ 延拓为 $G(\A)$ 上的函数：
\begin{equation*}\tag{1.10}
\phi(g)=\phi(\gamma g_\infty k)=\phi(g_\infty)\,\chi(k).
\end{equation*}
注意表达式 $g=\gamma g_\infty k$ 的含糊之处，只是把 $\gamma$ 换成 $\gamma\gamma_0^{-1}$、把 $g_\infty$ 换成 $\gamma_0g_\infty$、把 $k$ 换成 $\gamma_0k$，其中 $\gamma_0\in\Gamma_0(N)$。因此由引理 1.2 的 (i)，这个延拓是良定义的。注意 $\phi$ 仍满足该引理的 (ii)，并且在 $Z_\infty^+$（$Z_\infty$ 的单位连通分支）下不变。

至此我们还没有用到 $f$ 的全纯性。$f$ 的这一性质，等价于 $\phi$ 作为 $G(\mathbb{R})^+$ 上函数的如下性质。$G(\mathbb{R})^+$ 的实 Lie 代数 $\mathfrak{g}_0=\operatorname{Lie}(G(\mathbb{R}))=M_2(\mathbb{R})$ 通过
\[
X\phi(g_\infty)=\frac{d}{dt}\Bigl\{\phi\bigl(g_\infty\exp(tX)\bigr)\Bigr\}\Big|_{t=0}
\]
作用在 $G(\mathbb{R})^+$ 的光滑函数上。例如取
\[
X=\begin{pmatrix}0&1\\-1&0\end{pmatrix},
\]
则 $\exp(tX)=k_t\in K_\infty^+$；而引理 1.2 的 (ii)（它把 $f$ 的权编码在变换律 (1.1) 中）表明，对我们的 $\phi=\phi_f$ 有
\begin{equation*}\tag{1.11}
X\phi(g_\infty)=\frac{d}{dt}\Bigl\{\phi(g_\infty)(e^{it})^k\Bigr\}\Big|_{t=0}=ik\,\phi(g_\infty).
\end{equation*}
我们可以把 $\mathfrak{g}_0$ 的作用延拓为复化 $\mathfrak{g}=\mathfrak{g}_0\otimes_{\mathbb{R}}\mathbb{C}$ 的作用，方式是令 $(X+iY)\phi=X\phi+i\cdot Y\phi$。元素
\begin{equation*}\tag{1.12}
X_\pm=\frac12\begin{pmatrix}1&\pm i\\ \pm i&-1\end{pmatrix}\in\mathfrak{g}
\end{equation*}
具有如下性质：
\begin{equation*}\tag{1.13}
\operatorname{Ad}(k_\theta)X_\pm=k_\theta X_\pm k_\theta^{-1}=e^{\pm 2i\theta}X_\pm.
\end{equation*}
由此可知
\[
X_\pm\phi(g_\infty k_\theta)=\phi(g_\infty)\,(e^{i\theta})^{k\pm 2},
\]
所以 $X_\pm$ 把权升高或降低 $2$。把 $X_\pm\phi$ 经由 (1.9) 的权 $k\pm 2$ 版本对应于 $\mathbb{H}$ 上的函数 $M_\pm f$，并不难用原 $f$ 的导数表出。这样得到 $f$ 上的经典 Maass 算子 $M_\pm$，而方程 $M_-f=0$ 正是 Cauchy–Riemann 方程；例如见 [2, pp.\ 155, 129]。于是我们有下述结论。

\begin{annote}
原书编号在此处从 (1.13) 直接跳到 (1.15)，中间那段无编号的显示公式 $X_\pm\phi(g_\infty k_\theta)=\phi(g_\infty)(e^{i\theta})^{k\pm2}$ 实际上就是 (1.14)，排印时漏了号。译文照原书保留这一编号空缺，不代补：若擅自补上，(1.15) 及其后所有编号都会与原文错位，读者按号查对原书时会出错。引用本章公式请以页面位置为准，不宜只凭序号。这是 Birkhäuser 2004 年初版的排印遗漏，并非文字抽取造成的错误。

\par\vspace{0.15em}{\raggedleft 【来源：原书第 137 页（A）】\par}
\end{annote}

\medskip\noindent\textbf{引理 1.3.}\quad
满足 (1.1) 的光滑函数 $f:\mathbb{H}\to\mathbb{C}$ 全纯，当且仅当对应的函数 $\phi=\phi_f$ 满足 $X_-\phi=0$。

\medskip
至此，基本结论中除最后一部分外都已解释。

\medskip\noindent\textbf{命题 1.4.}\quad
映射 (1.8)、(1.9)、(1.10) 决定一个同构
\[
S_k(\Gamma_0(N),\chi_N)\overset{\sim}{\longrightarrow}\mathcal{A}_0(G)(\mathrm{hol},k,N,\chi),\qquad f\mapsto\phi_f,
\]
其中 $\mathcal{A}_0(G)(\mathrm{hol},k,N,\chi)$ 是 $G(\A)$ 上满足下述性质的函数 $\phi$ 构成的空间：
\begin{enumerate}[label=(\roman*),leftmargin=2.4em,itemsep=0.15em,topsep=0.25em]
\item 对所有 $\gamma\in G(\mathbb{Q})$ 与 $k\in K_0(N)$，$\phi(\gamma gk)=\phi(g)\chi(k)$。
\item 对任意 $g_f\in G(\A_f)$，函数 $g_\infty\mapsto\phi(g_\infty g_f)$ 在 $G(\mathbb{R})^+$ 上光滑、在 $Z_\infty^+$ 下不变，并且满足
\[
\phi(gk_\theta)=\phi(g)(e^{i\theta})^k
\quad\text{以及}\quad
X_-\phi=0.
\]
\item 函数 $\phi$ 是尖点的，即
\[
\int_{\mathbb{Q}\backslash\A}\phi\Bigl(\begin{pmatrix}1&x\\0&1\end{pmatrix}g\Bigr)\,dx=0
\]
对所有 $g\in G(\A)$ 成立。
\end{enumerate}

\medskip
剩下要解释的是 (iii)。对 $a\in\mathbb{R}^\times_+$，不难验证
\begin{equation*}\tag{1.15}
\begin{aligned}
\int_{\mathbb{Q}\backslash\A}\phi\Bigl(\begin{pmatrix}1&x\\0&1\end{pmatrix}\begin{pmatrix}a&\\&a^{-1}\end{pmatrix}\Bigr)dx
&=\int_{\mathbb{Z}\backslash\mathbb{R}}\phi\Bigl(\begin{pmatrix}1&x\\0&1\end{pmatrix}\begin{pmatrix}a&\\&a^{-1}\end{pmatrix}\Bigr)dx\\
&=a^k\int_0^1f(a^2i+x)\,dx=0,
\end{aligned}
\end{equation*}
因为 $f$ 在无穷远尖点处 Fourier 展开的常数项为零。另一方面，$f$ 的尖点性断言，类似的常数项在所有尖点处都为零。这等价于 (iii) 中的积分对一切 $g\in G(\A)$ 为零，也就是作为 $G(\A)$ 上的函数为零。

\begin{annote}
(1.15) 中间一步把积分限从 $\mathbb{Q}\backslash\mathbb{A}$ 换成 $\mathbb{Z}\backslash\mathbb{R}$，末行的 $a^k$ 也常被当作笔误，两处都需要核对。取 $g_\infty=\begin{pmatrix}a&\\&a^{-1}\end{pmatrix}\begin{pmatrix}1&u\\0&1\end{pmatrix}$，则 $g_\infty(i)=a^2i+a^2u$，$\det g_\infty=1$，故 $j(g_\infty,i)=a^{-1}$，(1.9) 给出 $\phi(g_\infty)=a^kf(a^2i+a^2u)$，$a^k$ 即由此而来。又
\[
\begin{pmatrix}1&x\\0&1\end{pmatrix}\begin{pmatrix}a&\\&a^{-1}\end{pmatrix}
=\begin{pmatrix}a&\\&a^{-1}\end{pmatrix}\begin{pmatrix}1&a^{-2}x\\0&1\end{pmatrix},
\]
令 $u=a^{-2}x$ 便把 $\mathbb{Z}\backslash\mathbb{R}$ 上的积分搬到 $a^{-2}\mathbb{Z}\backslash\mathbb{R}$ 上；$dx=a^2du$ 与周期长度 $a^{-2}$ 相消，于是积分等于 $a^k\int_0^1f(a^2i+y)\,dy$。这个积分是 $f$ 沿水平线的积分，等于其 Fourier 展开的常数项，故为零。至于积分限为何是 $\mathbb{Z}\backslash\mathbb{R}$：$\phi$ 对 $\begin{pmatrix}1&n\\0&1\end{pmatrix}$（$n\in\mathbb{Z}$）严格不变（引理 1.2(i) 中 $d=1$），故 $\mathbb{Q}\backslash\mathbb{A}=\mathbb{Z}\backslash\mathbb{R}\times\prod_p\mathbb{Z}_p$ 中 $\prod_p\mathbb{Z}_p$ 部分按归一化测度体积为 $1$，只剩 $\mathbb{Z}\backslash\mathbb{R}$。

\par\vspace{0.15em}{\raggedleft 【来源：由 (1.5)、(1.9) 与引理 1.2 直接算出；原书第 138 页（A）】\par}
\end{annote}

\medskip\noindent\textbf{注 1.5.}\quad
函数 $\phi$ 其实是中心 $Z(\A)$ 的本征函数。更确切地说，若 $\omega$ 是对应于 $\chi_N$ 或 $\chi$ 的 $\A^\times/\mathbb{Q}^\times$ 的特征，则
\begin{equation*}\tag{1.16}
\phi(zg)=\omega(z)\phi(g).
\end{equation*}

\medskip\noindent\textbf{注 1.6.}\quad
考虑函数
\begin{equation*}\tag{1.17}
\tilde\phi(g)=\phi(gI),\qquad I=\begin{pmatrix}1&\\&-1\end{pmatrix}\in O(2)=:K_\infty\subset G(\mathbb{R}).
\end{equation*}
注意
\begin{equation*}\tag{1.18}
\tilde\phi(gk_\theta)=\tilde\phi(g)\,(e^{i\theta})^{-k},
\end{equation*}
并且由于 $\operatorname{Ad}(I)X^\pm=X^\mp$，有
\begin{equation*}\tag{1.19}
X_+\tilde\phi=0.
\end{equation*}
不难验证，$\mathbb{H}$ 上对应的函数是
\begin{equation*}\tag{1.20}
\tilde f(\tau)=\tilde\phi(g_\infty)\,j(g_\infty,i)^{-k}=v^k f(-\bar\tau).
\end{equation*}

\section{$\mathcal{A}_0(G)$ 上的表示}

为了把群 $G(\A)$ 的表示论引进来，我们要扩大上一节构造的、$G(\mathbb{Q})Z_\infty^+\backslash G(\A)$ 上的函数空间 $\mathcal{A}_0(G)(\mathrm{hol},k,N,\chi)$。\footnote{人们起初也许会对 $G(\A)$ 中元素 $g$ 对 $\mathcal{A}_0(G)(\mathrm{hol},k,N,\chi)$ 中函数作任意右平移所得的组合加以尝试。结果发现，更好的做法是坚持要求函数都是“有权的”函数的有限线性组合，即都是 $K_\infty^+=\mathrm{SO}(2)$ 的本征函数。这样的函数空间在 $G(\mathbb{R})$ 的元素右平移下并不保持，因此只好退而求其次，满足于复化 Lie 代数 $\mathfrak{g}$ 的作用。} 设 $\omega$ 是 $\A^\times/\mathbb{Q}^\times$ 的特征，且 $\omega_\infty$ 在 $\mathbb{R}^\times_+$ 上平凡。

\medskip\noindent\textbf{定义 2.1.}\quad
$G(\A)$ 上中心特征为 $\omega$ 的自守形式空间 $\mathcal{A}(G)=\mathcal{A}(G,\omega)$，是 $G(\A)$ 上复值函数 $\phi$ 构成的空间，使得：
\begin{enumerate}[label=(\roman*),leftmargin=2.4em,itemsep=0.15em,topsep=0.25em]
\item 对 $z\in Z(\A)\simeq\A^\times$ 与 $\gamma\in G(\mathbb{Q})$，$\phi(z\gamma g)=\omega(z)\phi(g)$。
\item 对每个 $g_f\in G(\A_f)$，函数 $g_\infty\mapsto\phi(g_\infty g_f)$ 在 $G_\infty$ 上光滑。
\item 由 $\phi$ 经 $K_\infty$ 中元素右平移张成的空间是有限维的，即 $\phi$ 是右 $K_\infty$-有限的。
\item 存在 $G(\A_f)$ 的紧开子群\footnote{不妨设 $K'\subset K$。$K'$-不变性这一条件，等价于要求 $\phi$ 经 $K$ 中元素右平移张成的空间是有限维的，即 $\phi$ 是右 $K$-有限的。} $K'$，使得 $\phi$ 在右平移下 $K'$-不变。
\item 函数 $\phi$ 是 $\mathcal{Z}(\mathfrak{g})$-有限的。
\item 函数 $\phi$ 是缓增的。
\end{enumerate}

\medskip
\begin{annote}
定义 2.1 之前假设 $\omega_\infty$ 在 $\mathbb{R}^\times_+$ 上平凡，这与引理 1.2(ii) 的 $\operatorname{sgn}(z)^k$ 有一处不易察觉的张力。对奇权 $k$，由 (1.10) 构造的 $\phi_f$ 满足 $\phi(zg)=\operatorname{sgn}(z)^k\phi(g)$（因 $Z_\infty\subset G(\mathbb{R})^+$ 且 $\det z=z^2>0$），其 $\omega_\infty$ 是 $\operatorname{sgn}$ 而非平凡特征。故取 $\omega_\infty$ 平凡时，$\mathcal{A}(G,\omega)$ 只容纳偶权形式；奇权情形应改取 $\omega_\infty=\operatorname{sgn}^k$。好在本章后文只用到 $\omega_\infty$ 在 $\mathbb{R}^\times_+$ 上平凡（即 $\operatorname{Lie}(Z_\infty)$ 平凡作用）这一点，而 $\operatorname{sgn}$ 在 $\mathbb{R}^\times_+$ 上确实平凡，故两种取法下 (v)、定理 2.6 等结论均不变。译文照原文，不代改。

\par\vspace{0.15em}{\raggedleft 【来源：对照引理 1.2(ii) 与定义 2.1(i)，原书第 136、138--139 页（A）】\par}
\end{annote}

条件 (v) 的含义如下。由于光滑性条件 (ii)，复化 Lie 代数 $\mathfrak{g}$ 作用在 $\phi\in\mathcal{A}_0(G)$ 上，这一点已在第 1 节说明。注意这一作用保持条件 (i)–(iv)，特别是 $K_\infty$-有限性条件 (iii)。$\mathfrak{g}$ 的作用可延拓为 $\mathfrak{g}$ 的泛包络代数 $U(\mathfrak{g})$ 的作用。该代数的中心 $\mathcal{Z}(\mathfrak{g})$ 由 $\operatorname{Lie}(Z_\infty)=\operatorname{Lie}(Z_\infty^+)$ 与 Casimir 算子张成；前者由 (i) 及我们对 $\omega$ 的假设平凡地作用，而
\begin{equation*}\tag{2.1}
C=H^2+2X_+X_-+2X_-X_+,\qquad\text{其中}\quad
H=-i\begin{pmatrix}&1\\-1&\end{pmatrix},
\end{equation*}
$X_\pm$ 如 (1.12) 所示；参看 [2, Section 2.2]。于是条件 (v) 说的是：存在多项式 $R[T]\in\mathbb{C}[T]$，使得 $R(C)\phi=0$。

最后，为了表述增长条件 (vi)，人们使用 $G(\A)$ 上的下述范数。存在嵌入
\begin{equation*}\tag{2.2}
G(\A)\hookrightarrow M_2(\A)\times\A^\times\simeq\A^5,\qquad g\mapsto(g,\det(g)^{-1}).
\end{equation*}
对每个位 $v$，以及 $x=(x_i)\in\A^5$，置
\begin{equation*}\tag{2.3}
\|x\|_v=\max_i(|x_i|_v),
\end{equation*}
并令
\begin{equation*}\tag{2.4}
\|x\|=\prod_v\|x\|_v.
\end{equation*}
注意，例如当 $v$ 是非阿基米德位、$g_v\in K_v$ 时，$\det(g_v)$ 是单位，故 $\|g\|_v=1$。因此乘积中几乎所有的因子都等于 $1$。

于是，函数 $\phi$ 称为缓增的（或具适度增长的），如果存在整数 $n$ 与常数 $C$，使得
\begin{equation*}\tag{2.5}
|f(g)|\le C\,\|g\|^n
\end{equation*}
对所有 $g\in G(\A)$ 成立。

\begin{annote}
(2.5) 左端写作 $|f(g)|$，但此处 $g\in G(\A)$，而经典模形式 $f$ 只定义在上半平面 $\mathbb{H}$ 上（见 (1.2)）；本节此前已约定用 $\phi$ 记 $G(\A)$ 上的函数。故 (2.5) 中的 $f$ 应为 $\phi$。这是原书排印之误，同页上文皆为 $\phi$；它不影响缓增性的定义，因为增长条件本就是加在 $\phi$ 上的。译文照原文保留 $f$，以免擅自改动公式；引用时请按 $\phi$ 理解，否则会误以为这里对模形式本身加了范数条件。

\par\vspace{0.15em}{\raggedleft 【来源：原书第 139--140 页；Bump 上书第 2 章关于缓增性的定义（A）】\par}
\end{annote}

\medskip\noindent\textbf{定义 2.2.}\quad
中心特征为 $\omega$ 的尖点自守形式空间 $\mathcal{A}_0(G)=\mathcal{A}_0(G,\omega)\subset\mathcal{A}(G)$，由再加上下述尖点条件来定义：
\begin{enumerate}[label=(\roman*),leftmargin=2.4em,itemsep=0.15em,topsep=0.25em,start=7]
\item 函数 $\phi$ 是尖点的，即
\begin{equation*}\tag{2.6}
\int_{\mathbb{Q}\backslash\A}\phi\Bigl(\begin{pmatrix}1&x\\0&1\end{pmatrix}g\Bigr)\,dx=0
\end{equation*}
对所有 $g\in G(\A)$ 成立。
\end{enumerate}

\medskip
容易看出，空间 $\mathcal{A}_0(G)(\mathrm{hol},k,N,\chi)$ 含于 $\mathcal{A}_0(G)$。事实上，借助命题 1.4 以及注 1.5、注 1.6，(i)–(iv) 与 (vii) 都是直接的。命题 1.4 的 (ii) 中权与全纯性条件表明\footnote{由权条件 $H\phi=k\phi$。此外 $X_-\phi=0$ 且 $[X_+,X_-]=H$，故 $X_-X_+\phi=(X_+X_--H)\phi=-k\phi$。}
\begin{equation*}\tag{2.7}
C\phi=k(k-2)\phi.
\end{equation*}

\begin{annote}
(2.1) 的 Casimir 算子必须与 (1.12)、(2.9) 的归一化配套核对，否则 (2.7) 会差一个符号或因子。取 $H=-i\begin{pmatrix}1&\\&-1\end{pmatrix}$、$X_\pm=\frac12\begin{pmatrix}1&\pm i\\ \pm i&-1\end{pmatrix}$，直接计算得 $[X_+,X_-]=H$，于是 $X_-X_+\phi=(X_+X_--H)\phi=-k\phi$（此处用 $X_-\phi=0$ 与 $H\phi=k\phi$），这正是 (2.7) 后脚注中的算法。在 $\mathrm{DS}_k$ 上验算：由 $\pi_k(X_\pm)v_\ell=\frac12(k\pm\ell)v_{\ell\pm2}$ 得
\[
2X_+X_-v_\ell=\tfrac12(k-\ell)(k+\ell-2)v_\ell,\qquad
2X_-X_+v_\ell=\tfrac12(k+\ell)(k-\ell-2)v_\ell,
\]
与 $H^2v_\ell=\ell^2v_\ell$ 相加恰得 $(k^2-2k)v_\ell$。若把 $H$ 误读为 $\diag(1,-1)$，或漏掉 $2X_+X_-$ 与 $2X_-X_+$ 中对称的一项，都会算出别的特征值，并连带使定理 2.6 与第 4 节的对应关系错位。

\par\vspace{0.15em}{\raggedleft 【来源：原书第 139--140 页及其脚注；Bump 上书第 2 章（A）】\par}
\end{annote}

缓增性条件可由经典的尖点全纯性条件推出 [19]；而尖点函数在 $G(\A)$ 上其实是速降的 [2]。另一方面，第 1 节的构造可以施于经典（全纯）Eisenstein 级数，所得到的 $G(\A)$ 上的函数属于 $\mathcal{A}(G)$，从而是缓增的，但不属于 $\mathcal{A}_0(G)$。

\medskip\noindent\textbf{定义 2.3.}\quad
(i) 一个 $(\mathfrak{g},K_\infty)$-模，是指带有 $\mathfrak{g}$ 与 $K_\infty$ 作用的复向量空间 $V$，使得 $V$ 中所有向量都是 $K_\infty$-有限的，并且这两个作用相容，即
\begin{equation*}\tag{2.8}
k\cdot X\cdot v=\operatorname{Ad}(k)X\cdot k\cdot v,
\end{equation*}
且对 $X\in\operatorname{Lie}(K_\infty)$，
\begin{equation*}\tag{2.9}
X\cdot v=\frac{d}{dt}\Bigl(\exp(tX)\cdot v\Bigr)\Big|_{t=0}.
\end{equation*}
(ii) 一个 $(\mathfrak{g},K_\infty)\times G(\A_f)$-模，是指一个 $(\mathfrak{g},K_\infty)$-模，带有 $G(\A_f)$ 的光滑作用，且与 $(\mathfrak{g},K_\infty)$ 的作用交换。这里“光滑”是指：每个向量 $v\in V$ 都被某个紧开子群 $K'\subset G(\A_f)$ 固定。

\medskip
注意，$v\in V$ 的 $K_\infty$-有限性在 (2.9) 中用来保证曲线 $\exp(tX)\cdot v$ 落在一个有限维向量空间里。

这些定义的要点在于：空间 $\mathcal{A}(G)$ 与 $\mathcal{A}_0(G)$ 都带有
$(\mathfrak{g},K_\infty)\times G(\A_f)$-模结构，并且 $\mathcal{A}_0(G)(\mathrm{hol},k,N,\chi)$ 是 $\mathcal{A}_0(G)$ 的子空间。事实上，有一点需要验证，即缓增性条件 (vi) 在 $\mathfrak{g}$ 的作用下保持不变。关于这一点的漂亮讨论（它显示了有限性条件 (iii) 与 (v) 的非平凡作用），见 Bump [2, Section 2.2，特别是 Theorem 2.9.2]。

现在设 $k$ 是任意数域，$G=\GL(2)$ 视为 $k$ 上的代数群。对 $k$ 的任一位 $v$，$G_v=\GL_2(k_v)$；若 $v$ 非阿基米德，则 $K_v=\GL_2(\mathcal{O}_v)$。与前面一样，
\[
G(\A)={\prod_v}'G_v
\]
是关于子群 $K_v$ 的限制乘积，
\[
K=\prod_{v\ \text{非阿基米德}}K_v
\]
是 $G(\A_f)$ 的紧开子群，而 $G(k)\hookrightarrow G(\A)$ 是离散子群，等等。我们可以完全仿照定义 2.1（相应地，定义 2.2）来定义 $G(\A)$ 上的自守形式空间 $\mathcal{A}(G)$（相应地，尖点自守形式空间 $\mathcal{A}_0(G)$）。当然，此时
\begin{equation*}
G_\infty=\prod_{v\ \text{阿基米德}}G(k_v)\simeq
\underbrace{\GL_2(\mathbb{R})\times\cdots\times\GL_2(\mathbb{R})}_{r_1}
\times\underbrace{\GL_2(\mathbb{C})\times\cdots\times\GL_2(\mathbb{C})}_{r_2},
\end{equation*}
\begin{equation*}
K_\infty=\prod_{v\ \text{阿基米德}}K_v\simeq
\underbrace{O(2)\times\cdots\times O(2)}_{r_1}
\times\underbrace{U(2)\times\cdots\times U(2)}_{r_2},
\end{equation*}
而中心 $\mathcal{Z}(\mathfrak{g})\simeq\otimes_v\mathcal{Z}(\mathfrak{g}_v)$ 也更大。正如 Tate 论文以统一的方式处理所有数域的 Hecke 特征，Jacquet–Langlands 理论 [11] 对任意 $k$ 都适用于空间 $\mathcal{A}(G)\supset\mathcal{A}_0(G)$。

为了描述 $(\mathfrak{g},K_\infty)\times G(\A_f)$ 在空间 $\mathcal{A}_0(G)$ 上的表示，我们先回顾抽象表示论的若干事实。

\medskip\noindent\textbf{定义 2.4.}\quad
(i) $(\mathfrak{g},K_\infty)\times G(\A_f)$-模 $(\pi,V)$ 称为容许的，如果对 $K_\infty\times K$ 的每个不可约表示 $\sigma$，$\sigma$ 在 $V$ 中的重数有限。
(ii) 模 $(\pi,V)$ 称为不可约的，如果它没有在 $\mathfrak{g}$、$K_\infty$ 与 $G(\A_f)$ 作用下保持的真子空间。

\medskip
有类似的局部定义：对阿基米德位 $v$，$(\mathfrak{g}_v,K_v)$-模 $(\pi,V)$ 称为容许的，如果对 $K_v$ 的任一不可约表示 $\sigma_v$，$\sigma_v$ 在 $V$ 中的重数有限。这里 $\mathfrak{g}_v=\operatorname{Lie}(G_v)\otimes\mathbb{C}$。$(\mathfrak{g}_v,K_v)$-模称为不可约的，如果它没有在 $\mathfrak{g}_v$ 与 $K_v$ 下不变的真子空间。

对非阿基米德位 $v$，复向量空间 $V$ 上 $G_v$ 的表示 $(\pi,V)$ 称为光滑的，如果每个 $\phi\in V$ 都被某个开子群固定；称为容许的，如果对每个开子群 $K'$，$K'$-不动向量空间 $V^{K'}$ 都是有限维的；称为不可约的，如果它没有真的不变子空间。此外，不可约容许表示 $(\pi,V)$ 称为无分歧的或球表示，如果 $K_v$-不动向量空间非零。此时其实有 $\dim V^{K_v}=1$。

下述可因子化定理，代替了 $\GL_1$ 情形的一个简单事实：Hecke 特征 $\omega$ 可写成局部特征之积 $\omega=\otimes_v\omega_v$。

\medskip\noindent\textbf{定理 2.5.}\quad
(i) 设对 $k$ 的每个阿基米德（相应地，非阿基米德）位 $v$，$(\pi_v,V_v)$ 是不可约容许 $(\mathfrak{g}_v,K_v)$-模（相应地，$G_v$ 的表示）。再设对几乎所有的非阿基米德位，表示 $(\pi_v,V_v)$ 无分歧，并且给定了一个基向量\footnote{这意思是：$V$ 是空间 $V_S=\otimes_{v\in S}V_v$ 的并，即正向极限；当 $T\supset S$ 时，嵌入 $V_S\hookrightarrow V_T$ 由 $x_S\mapsto x_S\otimes(\otimes_{v\in T\setminus S}\xi_v^0)$ 给出。} $\xi_v^0\in V_v^{K_v}$。则关于向量 $\xi_v^0$ 的限制张量积
\[
V={\otimes_v}'V_v
\]
是不可约容许 $(\mathfrak{g},K_\infty)\times G(\A_f)$-模。
(ii) 反之，若 $(\pi,V)$ 是不可约容许 $(\mathfrak{g},K_\infty)\times G(\A_f)$-模，则存在如 (i) 的一族 $\{(\pi_v,V_v)\}_v$，以及 $(\mathfrak{g},K_\infty)\times G(\A_f)$-模的同构
\[
{\otimes_v}'V_v\overset{\sim}{\longrightarrow}V.
\]

\medskip
两种情形下习惯上都写作 $\pi=\otimes_v\pi_v$，尽管这一记号在情形 (ii) 中略去了那个重要的同构。详细的讨论见 [2, Sections 3.3 and 3.4]。

现在回到尖点模形式构成的 $(\mathfrak{g},K_\infty)\times G(\A_f)$-模。关于其结构的主要结果如下。回忆中心特征 $\omega=\otimes_v\omega_v$ 已经固定。

\medskip\noindent\textbf{定理 2.6.}\quad
(i) 空间 $\mathcal{A}_0(G)=\mathcal{A}_0(G,\omega)$ 是不可约容许 $(\mathfrak{g},K_\infty)\times G(\A_f)$-模的代数直和，
\[
\mathcal{A}_0(G)=\bigoplus_\pi\mathcal{A}_0(\pi).
\]
这里 $\mathcal{A}_0(\pi)$ 是同构于 $\pi$ 的直和项。
(ii)（重数一；Jacquet–Langlands [11]）若 $\pi_1$ 与 $\pi_2$ 是不可约容许 $(\mathfrak{g},K_\infty)\times G(\A_f)$-模，则
\[
\pi_1\simeq\pi_2\ \Longrightarrow\ \mathcal{A}_0(\pi_1)=\mathcal{A}_0(\pi_2).
\]
(iii)（强重数一；Miyake、Casselman [8, p.\ 89]）设 $\pi_1=\otimes_v\pi_{1,v}$、$\pi_2=\otimes_v\pi_{2,v}$ 如 (ii)，并且对 $k$ 的某个有限位集 $S$ 之外的所有 $v$ 都有 $\pi_{1,v}\simeq\pi_{2,v}$。则 $\pi_1\simeq\pi_2$，且 $\mathcal{A}_0(\pi_1)=\mathcal{A}_0(\pi_2)$。
(iv)（超强重数一；Ramakrishnan [17]）设 $\pi_{1,v}\simeq\pi_{2,v}$ 对 $k$ 的某个体密度小于 $\frac18$ 的位集之外的所有 $v$ 成立。则 $\pi_1\simeq\pi_2$，且 $\mathcal{A}_0(\pi_1)=\mathcal{A}_0(\pi_2)$。

\medskip
\begin{annote}
(2.6)(iv) 中的 $1/8$ 不是随手取的常数，它出自 Ramakrishnan 附于 Taylor 那篇文章之后的附录（原书 [17]）。经典强重数一（(ii)、(iii)）只允许两个表示在有限个位处不同；若按位集的自然密度衡量，(iv) 把例外集放宽到密度小于 $1/8$，$1/8$ 是该文密度估计的产物，并非已知的普适下界。证明的手段是 $\GL(2)$ 的对称平方 $L$ 函数（其自守性由 Gelbart–Jacquet 建立）与 Rankin–Selberg $L$ 函数：局部同构的位越多，这些 $L$ 函数在 $s=1$ 附近的极点信息越强，最终迫使 $\pi_1\simeq\pi_2$。本书只引用结论，未给证明。

\par\vspace{0.15em}{\raggedleft 【来源：Ramakrishnan，appendix to Taylor, \emph{Invent.\ Math.}\ \textbf{116}(1994), 645--649（原书 [17]）；Gelbart–Jacquet, \emph{Ann.\ Sci.\ ENS} \textbf{11}(1978)（A）】\par}
\end{annote}

出现在 $\mathcal{A}_0(G)$ 中的不可约容许 $(\mathfrak{g},K_\infty)\times G(\A_f)$-模 $\pi$，就是 $G(\A)$ 的不可约尖点自守表示。这样的表示称为在空间 $\mathcal{A}_0(\pi)$ 上实现。由于重数一，一旦抽象的表示 $\pi$ 给定，空间 $\mathcal{A}_0(\pi)$（若存在）就没有含糊之处。对其他经典群则不然。例如对 $G=\mathrm{SL}_n$，$n\ge 3$，给定的抽象表示 $\pi$ 可以在尖点形式空间中以大于一的重数出现 [1]。在这种情形下，指明空间 $\mathcal{A}_0(\pi)$ 可能很重要。断言 (i) 的证明基于两件事的结合：一方面，自守尖点表示 $\pi$ 中的函数 $\phi$ 由它的 Fourier 展开决定，而这些展开又由附于 $\phi$ 的整体 Whittaker 函数决定；另一方面，$\pi=\otimes_v\pi_v$ 中函数对应的这类 Whittaker 函数空间，由各局部成分 $\pi_v$ 对应的局部 Whittaker 函数空间唯一决定。细节可在 [11, 10, 8, 2] 中找到。

\section{局部表示论}

本节简要回顾 $G=\GL_2(F)$ 的表示论，其中 $F$ 是 $p$ 进域。本节末尾再补几句关于 $(\mathfrak{g},K_\infty)$-模的话。我们沿用 Tate 论文一章第 3 节的记号。于是 $F$ 是非阿基米德局部域，整数环为 $\mathcal{O}$。更详细的综述见 [8, Chapter 4]，详细证明例如见 [2, Chapter 4]。

取 $G=\GL_2(F)$，$K=\GL_2(\mathcal{O})$，于是 $K$ 是 $G$ 的极大紧（开）子群。设 $B=NM=MN$，其中
\[
M=\left\{m=m(a)=\begin{pmatrix}a_1&\\&a_2\end{pmatrix}\right\},
\qquad
N=\left\{n=n(x)=\begin{pmatrix}1&x\\&1\end{pmatrix}\right\},
\]
这里 $a_1,a_2\in F^\times$，$x\in F$。设 $Z$ 是 $G$ 的中心。

对 $F^\times$ 的特征对 $\mu=(\mu_1,\mu_2)$，有诱导表示
\[
I(\mu)=\Bigl\{f:G\to\mathbb{C}\ \Big|\ f(bg)=\mu(b)\delta(b)^{\frac12}f(g)\Bigr\},
\]
其中 $f$ 光滑，即在一个开子群 $K'\subset G$ 下右不变（$K'$ 可依赖于 $f$，不妨设它是 $K$ 的子群）；并且对 $b=nm(a)$，
\[
\mu(b)=\mu_1(a_1)\mu_2(a_2),\qquad \delta(b)=\left|\frac{a_1}{a_2}\right|.
\]
注意 $\delta$ 是 $B$ 在 $N$ 上伴随作用的模，即
\[
d(bnb^{-1})=\delta(b)\,dn
\]
对 $N$ 上的 Haar 测度 $dn$ 成立。因此，这里的诱导是归一化的，使得酉特征对 $(\mu_1,\mu_2)$ 给出可酉化的表示 $I(\mu)$。由于 Iwasawa 分解 $G=BK$，$I(\mu)$ 中的函数由它们在 $K$ 上的限制决定，从而
\[
I(\mu)\simeq I^K_{K\cap B}(\mu).
\]
又注意 $I(\mu)$ 的中心特征是 $\omega=\mu_1\mu_2$。

\medskip\noindent\textbf{定理 3.1.}\quad
\begin{enumerate}[label=(\roman*),leftmargin=2.4em,itemsep=0.15em,topsep=0.25em]
\item $I(\mu)$ 不可约，当且仅当 $\mu_1\mu_2^{-1}\ne\omega_{\pm1}$，其中 $\omega_s(x)=|x|^s$。
\item 若 $\mu_1\mu_2^{-1}=\omega_1$，则 $I(\mu)$ 有一个一维商，$G$ 在其上通过特征 $\chi\circ\det$ 作用，其中 $\mu=(\chi\omega_{\frac12},\chi\omega_{-\frac12})$；还有一个无限维不可约子表示 $\sigma(\mu)$。
\item 若 $\mu_1\mu_2^{-1}=\omega_{-1}$，则 $I(\mu)$ 有一个一维子模，$G$ 在其上通过特征 $\chi\circ\det$ 作用，其中 $\mu=(\chi\omega_{-\frac12},\chi\omega_{\frac12})$；还有一个无限维不可约商 $\sigma(\mu)$。
\item 这些表示之间仅有的等价关系如下。令 $\mu'=(\mu_2,\mu_1)$。若 $\mu_1\mu_2^{-1}\ne\omega_{\pm1}$，则 $I(\mu)\simeq I(\mu')$；若 $\mu_1\mu_2^{-1}=\omega_{\pm1}$，则 $\sigma(\mu)\simeq\sigma(\mu')$。
\end{enumerate}

\medskip\noindent\textbf{定理 3.2.}\quad
对 $G$ 的不可约容许表示 $(\pi,V)$，下述条件等价：
\begin{enumerate}[label=(\roman*),leftmargin=2.4em,itemsep=0.15em,topsep=0.25em]
\item 对所有 $\mu$，$\operatorname{Hom}_G(\pi,I(\mu))=0$；
\item $\pi$ 的矩阵系数\footnote{它们的定义如下。$V$ 的反特征表示 $\tilde V$ 是全对偶表示 $V^*=\operatorname{Hom}_{\mathbb{C}}(V,\mathbb{C})$ 中的光滑向量空间。对 $v\in V$ 与 $v'\in\tilde V$，$\phi_{v,v'}(g)=\langle\pi(g)v,v'\rangle$，其中 $\langle\,,\rangle$ 是 $(\pi,V)$ 与 $(\tilde\pi,\tilde V)$ 之间的配对。} $\phi_{v,v'}$ 模 $Z$ 紧支，即 $\phi_{v,v'}$ 的支集在 $G/Z$ 中的像是紧的。
\end{enumerate}

\medskip
\begin{annote}
定理 3.2(ii) 原文作 “compactly supported module $Z$”，其中 module 是旧式（法语）用法，意为“模”，此处应读作 modulo $Z$，即“模去中心 $Z$”。从数学上看也唯有这一读法：不可约容许表示的矩阵系数在中心方向上按中心特征变换，$Z$ 非紧，故其支集在 $G$ 中不可能紧（除非恒为零）；只有把 $\phi_{v,v'}$ 看作 $G/Z$ 上的函数，“紧支”才有意义，而“模中心紧支”正是超尖点表示的定义性特征。译文统一写作“模 $Z$ 紧支”，不保留原文的拼写。

\par\vspace{0.15em}{\raggedleft 【来源：原书第 144 页；Bump 上书第 4 章、Gelbart 上书第 4 章（A）】\par}
\end{annote}

满足定理 3.2 条件的表示，就是 $G$ 的超尖点表示。它们与不可约主序列表示 $I(\mu)$、特殊表示 $\sigma(\mu)$ 以及一维表示 $\chi\circ\det$ 一起，在同构意义下给出了 $G$ 的全部不可约容许表示。

回忆 $G$ 的不可约容许表示 $(\pi,V)$ 称为无分歧的，如果 $V^K\ne 0$，即有被极大紧子群 $K$ 固定的向量。若特征 $\mu=(\mu_1,\mu_2)$ 无分歧，即
\[
\mu_j(x)=t_j^{\operatorname{ord}(x)}
\]
对某对 $(t_1,t_2)\in(\mathbb{C}^\times)^2$ 成立，则 $\mu|_{B\cap K}=1$，而限制到 $K$ 给出同构
\[
I(\mu)\simeq C^\infty(B\cap K\backslash K).
\]
这个同构对 $K$ 的右乘作用是等变的。因此 $\dim I(\mu)^K=1$。

事实上，这一构造涵盖了所有无分歧表示。

\medskip\noindent\textbf{定理 3.3.}\quad
\begin{enumerate}[label=(\roman*),leftmargin=2.4em,itemsep=0.15em,topsep=0.25em]
\item 对每对 $(t_1,t_2)\in(\mathbb{C}^\times)^2$，存在 $G$ 的不可约容许无分歧表示 $\pi(t_1,t_2)$。
\item $G$ 的每个不可约容许无分歧表示 $\pi$ 都同构于某个 $\pi(t_1,t_2)$。这类表示之间唯一的等价关系是
\[
\pi(t_2,t_1)\simeq\pi(t_1,t_2).
\]
\item 若 $\mu_1\mu_2^{-1}\ne\omega_{\pm1}$，即 $t_1t_2^{-1}\ne q^{\pm1}$，则 $\pi(t_1,t_2)=I(\mu)$。若 $\mu_1\mu_2^{-1}=\omega_{\pm1}$，则 $\pi(t_1,t_2)=\chi\circ\det$，即 $I(\mu)$ 的一维成分。
\end{enumerate}

\medskip\noindent\textbf{注.}\quad
令
\[
I=\left\{k=\begin{pmatrix}a&b\\c&d\end{pmatrix}\in K\ \middle|\ \operatorname{ord}(c)\ge 1\right\}
\]
为 $K$ 的 Iwahori 子群。则 $G$ 的满足 $V^I\ne0$ 的不可约容许表示 $(\pi,V)$，只有无分歧表示与无分歧特殊表示，即 $\mu$ 无分歧时的 $\sigma(\mu)$。

无分歧表示的参数化可以表述如下。

\medskip\noindent\textbf{推论 3.4.}\quad
$G$ 的不可约容许无分歧表示的同构类，与复群 $G^\vee:=\GL_2(\mathbb{C})$ 中的半单共轭类之间存在双射，由
\[
\pi=\pi(t_1,t_2)\ \longleftrightarrow\ t(\pi):=\text{$\begin{pmatrix}t_1&\\&t_2\end{pmatrix}$ 的共轭类}
\]
给出。

\medskip
半单共轭类 $t(\pi)$ 称为 $\pi$ 的 Satake 参数。对任意连通约化群 $G$，Satake 证明的类似参数化 [18] 涉及 Langlands 对偶群 $^{L}G$。事实上，Satake 对 $G$ 的无分歧表示的描述，事后可视为 Langlands 对偶性的一个简单情形；它正是提示人们构造 $^{L}G$、并认识其在理解 $G$ 的表示论中的根本重要性的诸事之一 [15]。

\begin{annote}
推论 3.4 中的 $G^\vee=\GL_2(\mathbb{C})$ 是复对偶群，而非完整的 Langlands 对偶群 $^{L}G$。后者还须带上 Galois（或 Weil 群）部分：对分裂的 $\GL_2$ 有 $^{L}G=\GL_2(\mathbb{C})\times W_\mathbb{Q}$，其中 $W_\mathbb{Q}$ 的作用平凡。Satake 参数 $t(\pi)=\diag(t_1,t_2)$ 记录的是无分歧表示在球 Hecke 代数下的特征标 $\lambda_\pi$；而 $\mathcal{H}$ 同构于 $G^\vee$ 上半单共轭类空间的正则函数代数，于是“表示论侧的分类参数即对偶群中的半单共轭类”这条 Langlands 原则，在无分歧情形下成为一条定理。$t(\pi_v)$ 与第 5 节欧拉因子中所扮演的角色，恰对应 Galois 侧的 Frobenius 共轭类 $\mathrm{Fr}_v$；函子性猜想正是要在两类参数之间建立联系。

\par\vspace{0.15em}{\raggedleft 【来源：Satake, \emph{Publ.\ IHES} \textbf{18}(1963)（原书 [18]）；Langlands, \emph{Lect.\ Notes Math.}\ \textbf{170}(1970)（原书 [15]）；Bump 上书第 4 章（A）】\par}
\end{annote}

（球）Hecke 代数 $\mathcal{H}=\mathcal{H}(G,K)$，是 $G$ 上紧支且在左右乘法下都 $K$-不变的函数 $\xi$ 构成的空间，即 $\xi(k_1gk_2)=\xi(g)$。$\mathcal{H}$ 在卷积下成为代数，单位元是 $K$ 的特征函数\footnote{我们固定 $G$ 上的 Haar 测度，使 $K$ 的体积为 $1$。}。对 $G$ 的任一无分歧表示 $(\pi,V)$，$\mathcal{H}$ 通过
\[
\pi(\xi)v=\int_G\xi(g)\,\pi(g)v\,dg=\lambda_\pi(\xi)\,v
\]
作用在空间 $V^K$ 上，其中 $\lambda_\pi:\mathcal{H}\to\mathbb{C}$ 是代数同态。事实上，每个代数同态 $\mathcal{H}\to\mathbb{C}$ 都这样得到；并且经由映射 $\xi\mapsto\lambda_\pi(\xi)$，$\mathcal{H}$ 可等同于 $G^\vee$ 中半单共轭类空间上的正则函数代数。于是 $\mathcal{H}$ 有代数生成元 $\xi_1$ 与 $\xi_2$，经由代数同构
\[
\mathcal{H}=\mathbb{C}[\xi_1,\xi_2]\overset{\sim}{\longrightarrow}\mathbb{C}[t_1+t_2,t_1t_2],
\qquad
\begin{cases}
\xi_1\mapsto t_1+t_2=\lambda_\pi(\xi_1),\\[2pt]
\xi_2\mapsto t_1t_2=\lambda_\pi(\xi_2),
\end{cases}
\]
它们对应于标准对称多项式，其中 $\pi=\pi(t)$。

最后，下述结果对理解与经典理论的联系很重要。对整数 $r\ge0$，令
\[
K_0(r)=\left\{k=\begin{pmatrix}a&b\\c&d\end{pmatrix}\in K\ \middle|\ \operatorname{ord}(c)\ge r\right\}.
\]
若 $\omega$ 是 $F^\times$ 的特征、导子 $c(\omega)\le r$，即 $\omega$ 在 $1+\mathfrak{P}^r$ 上平凡，则 $\omega$ 通过
\[
\omega(k)=\omega(a)
\]
定义 $K_0(r)$ 的一个特征。

下述结果是 Casselman 的 [3]；参看 [8, Theorem 4.24]。

\medskip\noindent\textbf{定理 3.5.}\quad
对 $G$ 的任一不可约容许表示 $(\pi,V)$，中心特征为 $\omega_\pi$，存在最小整数 $c(\pi)\ge c(\omega_\pi)$，使得
\[
V^{K_0(r),\omega_\pi}:=\{v\in V\mid \pi(k)v=\omega_\pi(k)v,\ \forall k\in K_0(r)\}\ne0 .
\]
此外，
\[
\dim V^{K_0(c(\pi)),\omega_\pi}=1 .
\]

\medskip
整数 $c(\pi)$ 称为 $\pi$ 的导子。例如，无分歧表示有 $c(\pi)=0$，无分歧特殊表示有 $c(\pi)=1$，而所有其他无限维不可约容许 $\pi$ 都有 $c(\pi)\ge2$，特别是所有超尖点表示都如此。

\begin{annote}
定理 3.5 的 $c(\pi)$ 与第 4 节的 $N(\pi)=\prod_p p^{c(\pi_p)}$ 是同一件事的局部面与整体面：新形式的级恰等于各局部导子之积。$V^{K_0(r),\omega_\pi}$ 的非零性可借 Hecke 算子 $\xi_{p,1}$ 的特征值逐层检验，Casselman 定理断言“首次出现”的那一层恰为一维，这正是新形式唯一性的表示论根源，也与 Atkin–Lehner 的经典理论 [3] 相合。数值上：$c(\pi)=0$ 对应无分歧（$p\nmid N$），$c(\pi)=1$ 对应无分歧特殊表示（$p$ 恰整除级），只有 $c(\pi)\ge2$ 才可能是超尖点。故超尖点局部成分只出现在 $p^2\mid N$ 的位置，而 $p\|N$ 处的新形式在局部必是特殊表示。

\par\vspace{0.15em}{\raggedleft 【来源：Casselman, \emph{Math.\ Ann.}\ \textbf{201}(1973)（原书 [3]）；Gelbart 上书 Theorem 4.24；Atkin–Lehner, \emph{Math.\ Ann.}\ \textbf{185}(1970)（A）】\par}
\end{annote}

本节最后给出阿基米德理论的一个简短例子。设 $G=\GL_2(\mathbb{R})$，$\mathfrak{g}_0=\operatorname{Lie}(G)$，$\mathfrak{g}=\mathfrak{g}_0\otimes\mathbb{C}$，如前。令 $K_\infty=O(2)$。对整数 $k\ge2$，存在不可约 $(\mathfrak{g},K_\infty)$-模 $\mathrm{DS}_k=(\pi_k,V)$，定义如下。作为复向量空间，
\[
V=\bigoplus_{\substack{\ell\ge k\\ \ell\equiv k\,(2)}}\mathbb{C}v_\ell ,
\]
其中 $K_\infty$ 的作用由
\[
\pi_k(k_\theta)v_\ell=e^{i\ell\theta}v_\ell
\quad\text{与}\quad
\pi_k\Bigl(\begin{pmatrix}1&\\&-1\end{pmatrix}\Bigr)v_\ell=v_{-\ell}
\]
给出。于是由 (2.9) 有 $\pi_k(H)v_\ell=\ell v_\ell$。此外，升降算子 (1.12) 的作用为
\[
\pi_k(X_\pm)v_\ell=\frac12(k\pm\ell)\,v_{\ell\pm2},
\]
从而 $\pi_k(X_-)v_k=0$，$\pi_k(X_+)v_{-k}=0$，且 $\pi_k(C)=k(k-2)$。容易验证这些公式定义了不可约 $(\mathfrak{g},K_\infty)$-模，即 $\GL_2(\mathbb{R})$ 的权为 $k$ 的离散序列表示的代数版本。这就为实位处不可约容许 $(\mathfrak{g},K_\infty)\times G(\A_f)$-模的阿基米德成分提供了一种可能。完整的列表例如见 [2, Theorem 2.5.2, p.\ 219]。

\section{经典新形式}

现在回到 $k=\mathbb{Q}$ 的情形，以及命题 1.4 的同构
\begin{equation*}\tag{4.1}
S_k(\Gamma_0(N),\chi_N)\overset{\sim}{\longrightarrow}\mathcal{A}_0(G)(\mathrm{hol},k,N,\chi),\qquad f\mapsto\phi_f .
\end{equation*}
若 $p\nmid N$，则右端空间中的函数在 $K_p$ 的右平移下不变，于是有局部 Hecke 代数 $\mathcal{H}_p=\mathcal{H}(G_p,K_p)$ 的作用。当转移到 $S_k(\Gamma_0(N),\chi_N)$ 上时，$\mathcal{H}_p$ 的生成元 $\xi_{p,1}$ 与 $\xi_{p,2}$ 的作用与经典 Hecke 算子 $T_p$ 和 $R_p$ 的关系为
\begin{equation*}\tag{4.2}
f_0\mid\xi_{p,1}=p^{-\frac{k-1}{2}}f_0\mid T_p
\quad\text{与}\quad
f_0\mid\xi_{p,2}=f_0\mid R_p ,
\end{equation*}
见 [8, p.\ 47]。因此，代数
\begin{equation*}\tag{4.3}
\mathcal{H}^N={\otimes_{p\nmid N}}'\mathcal{H}_p
\end{equation*}
的作用，等价于不整除级的素数所对应的经典 Hecke 算子 $T_n$ 的作用。这就使人们可以用表示论的语言表述经典的新形式与老形式理论。

先设 $\pi=\pi_\infty\otimes(\otimes_p\pi_p)$ 是不可约尖点自守表示，且 $\pi_\infty$ 同构于第 3 节末尾所述的离散序列表示 $\mathrm{DS}_k$。令 $\omega=\omega_\pi$ 是 $\pi$ 的中心特征，并令
\begin{equation*}\tag{4.4}
N=N(\pi)=\prod_p p^{c(\pi_p)},
\end{equation*}
其中 $c(\pi_p)$ 是局部成分 $\pi_p$ 的导子。设 $K_0(N)$ 如 (1.3)。则由第 3 节概述的结果，空间
\begin{equation*}\tag{4.5}
\begin{aligned}
\mathcal{A}_0(\pi)^{\mathrm{new}}_{\mathrm{hol}}
:=\Bigl\{\phi\in\mathcal{A}_0(\pi)\ \Big|\ {}
&\text{(i) } X_-\phi=0,\ \text{且}\\
&\text{(ii) } \phi(gk)=\omega(k)\,\phi(g),\ \forall k\in K_0(N(\pi))\Bigr\}
=\mathbb{C}\phi_0
\end{aligned}
\end{equation*}
的维数为 $1$，并且落在子空间 $\mathcal{A}_0(G)(\mathrm{hol},k,N,\chi)$ 中，其中 $\chi$ 是联系于 $\omega$ 的 Dirichlet 特征。由命题 1.4，存在唯一一个对应于 $\phi_0$ 的全纯尖点形式 $f_0\in S_k(\Gamma_0(N),\chi)$。

若 $p\nmid N$，则局部成分 $\pi_p$ 无分歧，局部 Hecke 代数 $\mathcal{H}_p=\mathcal{H}(G_p,K_p)$ 通过特征 $\lambda_{\pi_p}$ 作用在空间 $\pi_p^{K_p}$ 上。于是尖点形式 $f_0$ 是所有这类 $p$ 的 Hecke 算子的本征函数：
\begin{equation*}\tag{4.6}
f_0\mid T_p=a_pf_0=p^{\frac{k-1}{2}}(t_{p,1}+t_{p,2})f_0
\quad\text{与}\quad
f_0\mid R_p=\chi(p)f_0=t_{p,1}t_{p,2}f_0 ,
\end{equation*}
其中 $t_{p,1}$ 与 $t_{p,2}$ 是 $\pi_p$ 的 Satake 参数 $t(\pi_p)$ 的两个分量。注意，由于 $\omega_\pi(p)=t_{p,1}t_{p,2}$，这些关系等价于
\begin{equation*}\tag{4.7}
1-a_p\,p^{-\frac{k-1}{2}}X+\chi(p)X^2=(1-t_{p,1}X)(1-t_{p,2}X)=\det\bigl(\mathbf{1}_2-t(\pi_p)X\bigr),
\end{equation*}
其中 $a_p$ 是 $f$ 的 $q$ 展开 (1.2) 中的第 $p$ 个系数。于是，$\pi_p$ 的 Satake 参数可以从经典 Hecke 多项式即 (4.7) 的左端读出，反过来也一样。

\begin{annote}
(4.7) 的左端是经典 Hecke 多项式，右端是 $\det(\mathbf{1}_2-t(\pi_p)X)$，两者相等正是“局部因子＝特征多项式”这一原则的精确翻译。由 (4.7) 与 $t_{p,1}t_{p,2}=\chi(p)$，两个 Satake 参数被 $a_p$ 与 $\chi(p)$ 完全决定。Ramanujan–Petersson 猜想（Deligne 1974 年由 Weil 猜想推出）断言 $|a_p|\le2p^{(k-1)/2}$，这等价于 $|t_{p,1}|=|t_{p,2}|=p^{1/2}$，此时 $\pi_p$ 是主序列表示、$I(\mu)$ 可酉化。$p\mid N$ 时没有这样的界：$p\|N$ 时两参数的绝对值分别为 $p^{1/2}$ 与 $p^{-1/2}$，这种不对称正是特殊表示的标志，也正是 Ramanujan 猜想在坏位处必然失效的原因。

\par\vspace{0.15em}{\raggedleft 【来源：Deligne, \emph{Publ.\ IHES} \textbf{43}(1974)；Iwaniec–Kowalski《Analytic Number Theory》第 14 章（A）】\par}
\end{annote}

回忆：若 $f$ 是所有 $T_p$ 与 $R_p$（$(p,N)=1$）的本征函数，并且 $f$ 不“来自更低的级”，则 $f$ 是新形式，[7, 14, 16]。用 $S_k(\Gamma_0(N),\chi)^{\mathrm{new}}$ 记新形式张成的空间。

\medskip\noindent\textbf{定理 4.1.}\quad
(i) 联系于自守尖点表示 $\pi$ 的全纯尖点形式 $f_0=f_0(\pi)\in S_k(\Gamma_0(N),\chi)$ 是新形式。
(ii) 设 $\omega$ 是对应于 $\chi$ 的 Hecke 特征。则在命题 1.4 的同构下，
\[
S_k(\Gamma_0(N),\chi)^{\mathrm{new}}\cong
\bigoplus_{\substack{\pi\\ \omega_\pi=\omega\\ \pi_\infty\simeq\mathrm{DS}_k\\ N(\pi)=N}}
\mathcal{A}_0(\pi)^{\mathrm{new}}_{\mathrm{hol}},
\]
从而每个新形式在相差一个纯量的意义下都是某个 $f_0(\pi)$。
(iii) 此外，
\[
S_k(\Gamma_0(N),\chi)\cong\mathcal{A}_0(G)(\mathrm{hol},k,N,\chi)=
\bigoplus_{\substack{\pi\\ \omega_\pi=\omega\\ \pi_\infty\simeq\mathrm{DS}_k\\ N(\pi)\mid N}}
\mathcal{A}_0(\pi)^{K_0(N),\omega}_{\mathrm{hol}} .
\]

\medskip
这里 $\mathcal{A}_0(\pi)^{K_0(N),\omega}_{\mathrm{hol}}$ 是把 $\mathcal{A}_0(\pi)^{\mathrm{new}}_{\mathrm{hol}}$ 定义中条件 (ii) 里的 $N(\pi)$ 换成 $N$ 而得到的。

\medskip\noindent\textbf{注记.}\quad
(i) 这一结果给出了经典新形式理论与自守表示之间的词典。事实上，读者若愿意，可以把这一结果当作经典新形式、老形式等等的定义。
(ii) 注意 Hecke 算子的作用只涉及 $G(\A_f)$ 的表示，因此若把离散序列表示 $\mathrm{DS}_k$ 换成别的不可约容许 $(\mathfrak{g},K_\infty)$-模，它的表现方式完全相同。这就解释了给定权的 Maass 形式的 Hecke 理论，此时 $\pi_\infty$ 是合适的主序列表示。

最后，我们回到最初的出发点，即经典全纯尖点形式 $f\in S_k(\Gamma_0(N),\chi_N)$。对应的函数 $\phi\in\mathcal{A}_0(G)(\mathrm{hol},k,N,\chi)$ 可以按定理 4.1 的 (iii) 分解为若干成分。这相当于经典理论中把 $f$ 分解为新形式与老形式的组合。另一方面，若取由 $\phi$ 生成的 $(\mathfrak{g},K_\infty)\times G(\A_f)$-子模 $\mathcal{V}(\phi)\subset\mathcal{A}_0(G)$，则得到的恰是那些使 $\phi$ 在 (iii) 中对应成分非零的 $\mathcal{A}_0(\pi)$ 之和。若 $f$ 是新形式，则空间 $\mathcal{V}(\phi)$ 不可约。

\section{自守尖点表示的 $L$ 函数}

本节描述联系于数域 $k$ 上不可约尖点自守表示 $\pi=\otimes_v\pi_v$ 的 Langlands $L$ 函数。

首先回忆：对每个这样的 $\pi$，存在 $k$ 的有限位集 $S$（含阿基米德位），使得当 $v\notin S$ 时 $\pi_v$ 无分歧。于是由推论 3.4，存在 $G^\vee=\GL_2(\mathbb{C})$ 中的一族 Satake 参数 $\{t(\pi_v)\}_{v\notin S}$，它们对应于 $\pi$ 的无分歧局部成分。这一情形类似于：给定有限 Galois 扩张 $\mathcal{K}/k$，对每个在 $\mathcal{K}/k$ 中无分歧的有限位 $v$，有 $\operatorname{Gal}(\mathcal{K}/k)$ 中一个 Frobenius 元的共轭类 $\mathrm{Fr}_v$，从而有共轭类族 $\{\mathrm{Fr}_v\}_{v\notin S}$。仿照 Artin $L$ 函数，对 $\GL_2(\mathbb{C})$ 的任一有限维表示 $r$ [13]，以及 $s\in\mathbb{C}$，相应的（部分）Langlands $L$ 函数定义为如下欧拉乘积
\[
L^S(s,\pi,r):=\prod_{v\notin S}\det\bigl(1-r(t(\pi_v))q_v^{-s}\bigr)^{-1}.
\]
可以证明它在某个右半平面收敛；而本理论的基本问题之一是：(i) 通过补充 $v\in S$ 处的因子来完备化 $L$ 函数；(ii) 建立完备化函数的解析性质，即解析延拓与函数方程。对 $r$ 为标准表示的情形，Jacquet 与 Langlands [11] 建立了完整而漂亮的理论。

\begin{annote}
本节开头那个欧拉乘积只含 $v\notin S$ 的因子，所缺的正是坏位因子与阿基米德因子（伽马因子）。完备化后的 $\Lambda(s,\pi,r)$ 应满足函数方程 $\Lambda(s,\pi,r)=\varepsilon(s,\pi,r)\Lambda(1-s,\tilde\pi,r)$，其中 $\tilde\pi$ 为反特征表示，$\varepsilon$ 因子是局部 $\varepsilon$ 因子之积。对 $r$ 为标准表示，这就是 Jacquet–Langlands 的理论 [11]；对一般的 $r$，Langlands 函子性预测 $L(s,\pi,r)$ 仍是某个自守表示的 $L$ 函数，从而同样具备解析延拓与函数方程。还须留意归一化：欧拉因子用 $q_v^{-s}$（$q_v$ 为剩余域大小），而本节末行的 $L^N(s+\frac{k-1}{2},f)$ 表明经典 $L$ 函数与表示论 $L$ 函数相差一个中心平移 $\frac{k-1}{2}$，这是全章最容易读错的一处记号。

\par\vspace{0.15em}{\raggedleft 【来源：Jacquet–Langlands, \emph{Lect.\ Notes Math.}\ \textbf{114}(1970)（原书 [11]）；Cogdell，本卷 pp.\ 197--228（原书 [4]）（A）】\par}
\end{annote}

例如，在 Hecke 考虑的情形，$k=\mathbb{Q}$，$\pi$ 是对应于权为 $k$、级为 $N$ 的全纯新形式 $f$ 的不可约尖点自守表示，$f$ 的 $q$ 展开 (1.2) 已归一化使 $a_1=1$。取 $r$ 为 $\GL_2(\mathbb{C})$ 的标准二维表示，则由 (4.7)，
\begin{equation*}
\begin{aligned}
L^S(s,\pi,r)&=\prod_{p\nmid N}\bigl(1-a_p\,p^{-s-\frac{k-1}{2}}+p^{-2s}\bigr)^{-1}\\
&=L^N\Bigl(s+\frac{k-1}{2},f\Bigr)\\
&=\sum_{\substack{n\\ (n,N)=1}}a_n\,n^{-s}.
\end{aligned}
\end{equation*}
关于这类 $L$ 函数的出色而完整的讨论，读者可参阅 [2, 8, 9, 12, 7] 等。

\section*{参考文献}
\addcontentsline{toc}{section}{参考文献}
\markright{参考文献}

\begin{enumerate}[label={[\arabic*]},leftmargin=2.8em,itemsep=0.35em,topsep=0.4em]
\item D. Blasius, On multiplicities for $SL(n)$, \emph{Israel J.\ Math.}, \textbf{88}(1994), 1--3.
\item D. Bump, \emph{Automorphic Forms and Representations}, Cambridge University Press, Cambridge, UK, 1997.
\item W. Casselman, On some results of Atkin and Lehner, \emph{Math.\ Ann.}, \textbf{201}(1973), 301--314.
\item J. Cogdell, Analytic theory of $L$-functions for $\GL_n$, in J. Bernstein and S. Gelbart, eds., \emph{An Introduction to the Langlands Program}, Birkhäuser Boston, Boston, 2003, 197--228 (this volume).
\item J. Cogdell, Langlands conjectures for $\GL_n$, in J. Bernstein and S. Gelbart, eds., \emph{An Introduction to the Langlands Program}, Birkhäuser Boston, Boston, 2003, 229--249 (this volume).
\item J. Cogdell, Dual groups and Langlands functoriality, in J. Bernstein and S. Gelbart, eds., \emph{An Introduction to the Langlands Program}, Birkhäuser Boston, Boston, 2003, 251--269 (this volume).
\item F. Diamond and J. Im, Modular forms and modular curves, in \emph{Seminar on Fermat's Last Theorem}, CMS Conference Proceedings 17, American Mathematical Society, Providence, RI, 1995, 39--133.
\item S. Gelbart, \emph{Automorphic Forms on Adele Groups}, Princeton University Press, Princeton, NJ, 1975.
\item S. Gelbart, An elementary introduction to the Langlands program, \emph{Bull.\ Amer.\ Math.\ Soc.}, \textbf{10}(1984), 177--219.
\item R. Godement, \emph{Notes on the Jacquet--Langlands Theory}, Institute for Advanced Study, Princeton, NJ, 1970.
\item H. Jacquet and R. P. Langlands, \emph{Automorphic forms on $\GL(2)$}, Lecture Notes in Mathematics 114, Springer-Verlag, New York, 1970.
\item A. W. Knapp, Introduction to the Langlands program, \emph{Proc.\ Sympos.\ Pure Math.}, \textbf{61}(1997), 245--302.
\item E. Kowalski, Elementary theory of $L$-functions I and II, in J. Bernstein and S. Gelbart, eds., \emph{An Introduction to the Langlands Program}, Birkhäuser Boston, Boston, 2003, 1--37 (this volume).
\item S. Lang, \emph{Introduction to modular forms}, Springer-Verlag, New York, 2001.
\item R. P. Langlands, Problems in the theory of automorphic forms, in \emph{Lectures in Modern Analysis and Applications}, Lecture Notes in Mathematics 170, Springer-Verlag, New York, 1970, 18--86.
\item T. Miyake, \emph{Modular Forms}, Springer-Verlag, New York, 1989.
\item D. Ramakrishnan, A refinement of the strong multiplicity one theorem for $\GL(2)$, appendix to R. Taylor, “$l$-adic representations associated to modular forms over imaginary quadratic fields II,” \emph{Invent.\ Math.}, \textbf{116}(1994), 645--649.
\item I. Satake, Theory of spherical functions on reductive algebraic groups over $p$-adic fields, \emph{Publ.\ IHES}, \textbf{18}(1963), 1--69.
\item G. Shimura, \emph{Introduction to the Arithmetic Theory of Automorphic Forms}, Princeton University Press, Princeton, NJ, 1971.
\end{enumerate}
